% ==================== CHAPTER 4 ====================
\chapter{IMPLEMENTATION, RESULTS, AND DISCUSSION}
\clearpage

\section{Architectural Performance Evaluation}
During implementation, the proposed architecture was evaluated against a traditional Python/Spark/Elasticsearch stack. The empirical results validated the architectural choices and demonstrated profound improvements in high-concurrency environments (Industry Alignment).

\begin{table}[htbp]
\centering
\caption{Architecture Comparison: Proposed Stack (Go/Gorse) vs.\ Traditional (Python)}
\label{tab:arch-compare}
\begin{tabularx}{\textwidth}{XXXX}
\toprule
\textbf{\textit{Feature}} & \textbf{\textit{Proposed (Go/Gorse)}} & \textbf{\textit{Traditional (Python/Flask)}} & \textbf{\textit{Architectural Advantage}} \\
\midrule
\textbf{\textit{Concurrency}} & Go (Gin / Goroutines) & Python (Flask / WSGI) & Go natively handles thousands of simultaneous feed requests without blocking operations. \\
\rowcolor{gray!10}
\textbf{\textit{Model Selection}} & Automated (Gorse AutoML) & Manual Grid Search & Gorse dynamically optimizes tuning, transitioning between MF and FM as sparsity increases. \\
\textbf{\textit{Search Speed}} & Typesense (In-Memory Vector) & Elasticsearch (Disk Lexical) & Typesense guarantees consistent, sub-millisecond candidate generation for vectors. \\
\bottomrule
\end{tabularx}
\end{table}

\begin{table}[htbp]
\centering
\caption{Evaluation of System Latency (Sub-50ms) During High Concurrency}
\label{tab:latency}
\begin{tabularx}{\textwidth}{XXXX}
\toprule
\textbf{\textit{System Component}} & \textbf{\textit{Average Latency (Under High Load)}} & \textbf{\textit{Target Threshold}} & \textbf{\textit{Status}} \\
\midrule
\textbf{\textit{API Gateway (Gin)}} & $\sim$5 ms & $<$ 10 ms & Optimal \\
\rowcolor{gray!10}
\textbf{\textit{Vector Search (Typesense)}} & $\sim$18 ms & $<$ 50 ms & Optimal \\
\textbf{\textit{Ranking Core (Gorse)}} & $\sim$22 ms & $<$ 50 ms & Optimal \\
\bottomrule
\end{tabularx}
\end{table}

\section{User Experience (UX) and State Management}
A major achievement of the implementation was the frontend responsiveness. By leveraging Redux Toolkit (RTK) and RTK Query, an ``optimistic updates'' pattern was implemented. When a user interacts with a post, the UI updates instantaneously without waiting for the server's \texttt{200 OK} response. This masks network latency, ensuring interactions appear instantaneous---a critical requirement for modern social applications.

\section{Discussion and Defense of Stack (Industry Alignment)}
As part of the project defense, the rationale for this stack is grounded in industry alignment. Go and Gorse were selected to match the high-concurrency demands of real-time social media. Typesense was implemented specifically to guarantee sub-50ms candidate generation latency. The integration of RTK Query on the frontend ensures the seamless, optimistic user experience expected by today's social media consumers.

\section{Multi-Stakeholder Value and Algorithmic Diversity}
Following the modern evaluation paradigms established in the literature \cite{ricci2022}, the system's success is uniquely evaluated far beyond standard mathematical accuracy metrics such as MSE or point-wise precision. Relying solely on historical interaction accuracy inherently suppresses serendipity and punishes new or out-of-network content creators (defined as the ``Provider'' stakeholders). By implementing the ``Explore vs.\ Exploit'' logic within the final ranking phase, the architecture intrinsically guarantees \textit{novelty} and \textit{diversity} inside the feed. This ensures the end-user (the ``Consumer'') avoids psychological echo chambers or filter bubbles, while also providing emerging content creators with fair, unbiased exposure opportunities. Structuring the system through this multi-stakeholder paradigm resolves the core algorithmic paradox common to equitable social media platforms.
